\section{Model-field architecture: selection, witness tying, and the open norm problem}

\subsection{The finite-field selector is exact}

\begin{theorem}[PROVED: L4, the small-field exhaustion]
Let $q$ be an odd prime power with $q\leq 25$, let $\chi$ be the quadratic
character of $\F_q$ with $\chi(0)=0$, and let
$\tau\in\F_q\setminus\{0,\pm1\}$.  There is an $a\in\F_q$ such that
\[
 \chi(1+4a^2)=-1,
 \qquad
 \chi\bigl(-(1+4a^2)\,(1-\tau^2(1+4a^2))\bigr)=1.
\]
% src: THEOREMS.md:278-289
For $q=3$ the assertion is vacuous, because there is no admissible $\tau$.
For $q\in\{5,7,9,11,13\}$ the minimum number of choices is $2$, and for
$q\in\{17,19,23,25\}$ it is $4$.
% src: THEOREMS.md:287-294
\end{theorem}

Thus the constant $25$ in the unrefereed character-sum proof of
\cite{sun} is not a genuine small-field obstruction: among odd residue
fields, only the field of order $3$ remains, and there the issue is
admissibility rather than failure of the selector.
% src: THEOREMS.md:300-304

\begin{theorem}[PROVED: L5, the Legendre identity]
For every odd prime power $q$, every
$\tau\in\F_q\setminus\{0,\pm1\}$, and $\lambda=\tau^{-2}$, let
$E_\lambda:y^2=x(x-1)(x-\lambda)$.  Then
\[
N(\tau) := \#\bigl\{a \in \mathbb{F}_q : \chi(1+4a^2) = -1,\
\chi\bigl(-(1+4a^2)(1-\tau^2(1+4a^2))\bigr) = 1\bigr\} \;=\; \frac{\#E_\lambda(\mathbb{F}_q)}{4}.
\]
% src: THEOREMS.md:308-315
Consequently
\[
 \#E_\lambda(\F_q)\geq(\sqrt q-1)^2>0,
 \qquad N(\tau)\geq(\sqrt q-1)^2/4>0,
\]
so the hypothesis $|k|>25$ can be deleted from the lemma.
% src: THEOREMS.md:316-320
\end{theorem}

The proof is an exact point count, not an asymptotic substitution.  With
$u_a=1+4a^2$ and $w_a=-u_a(1-\tau^2u_a)$, expand the corrected indicator;
the hyperbola identity $\sum_b\chi(b^2+c)=-1$ for $c\neq0$ evaluates the
quadratic pieces.  Fibre-counting the map $a\mapsto u_a$ then gives
\[
4N=q+1+\sum_u\chi\bigl(u(u-1)(u-\lambda)\bigr)
   =\#E_\lambda(\F_q).
\]
% src: THEOREMS.md:322-359
This also explains divisibility by $4$: the Legendre curve has its full
rational $2$-torsion.
% src: THEOREMS.md:366-369

\subsection{The tied model and its six-variable bookkeeping}

Set $K=\Q$, $S=\{2\}$, $\pi=2$, and $u=1$.  Pull the bridge back along
$z\mapsto z^3$ and put
$A=1+4a^2$, $B=2b$, $\delta_\tau=1-A\tau^2$, and
$c=h(a,b,z^3)=a^2z^6g(a,b)/(1-z^3-a^2z^6)$.
% src: THEOREMS.md:413-420
The original block is
\[
\Psi_\tau(a,b,c):\ \exists y,r,s\ [\delta_\tau(c^2 - Ay^2) - 16B(r^2 - As^2) = 16].
\]
% src: THEOREMS.md:419-420
The architectural step is to identify its witness $s$ with the congruence
parameter by imposing $a=1+2s$.  We use the finite branch menu
\[
 \tau\in\left\{0,\frac{2a}{A},
 \tau^\dagger=\frac{1+2a^2}{A}\right\}.
\]
After clearing the nonzero denominators, the resulting two-witness block is
\[
\begin{aligned}
\Theta_*(a,b,c,s):\quad \exists y,r\quad&
\bigl[c^2-Ay^2-16Br^2=16-16ABs^2\bigr]\\
&{}\vee
\bigl[c^2-Ay^2-16ABr^2=16A-16A^2Bs^2\bigr]\\
&{}\vee
\bigl[a^4(c^2-Ay^2)+4ABr^2=4A^2Bs^2-4A\bigr].
\end{aligned}
\]
The three equations correspond respectively to
$\tau=0$, $\tau=2a/A$, and $\tau^\dagger$; on the last branch
$\delta_{\tau^\dagger}=-4a^4/A$ and
$\alpha=(2a^2)^2$ is a square.  A finite disjunction is represented by the
product of these polynomials in the same witnesses $(y,r)$, so adding the
square branch costs no variable.
% src: data/l19_classexist.jsonl (square-branch definition); THEOREMS.md:1100-1121
The candidate model is therefore
\[
F(z):=\exists b,s\ \bigl[(1+2s,b)\in\Phi_1^{\{2\}}\ \wedge\
\Theta_*(1+2s,b,h(1+2s,b,z^3),s)\bigr].
\]
% src: THEOREMS.md:433-439
Here $\Phi_1^{\{2\}}$ enforces $v_2(s)\geq0$ and
$b\in\mathcal O_2^\times$.  Membership in it has existential rank at most
$3$; the tied block has rank $2$ in $(y,r)$.
% src: THEOREMS.md:453-456,470-473
The intersection theorem of \cite{ddf} replaces their conjunction by
$3+2-1=4$ witnesses, and projection of the two base coordinates $(b,s)$
gives
\[
 2+(3+2-1)=6.
\]
% src: THEOREMS.md:453-465
If an input set has rank zero, the elementary conjunction bound is no
larger.  The two branches share $(y,r)$ and are combined by multiplying
their equations; denominator clearing introduces no inverse witness because
$A\equiv5\pmod8$ at $2$ and $1-z^3-a^2z^6$ cannot vanish over $\Q$.
% src: THEOREMS.md:441-463

For use by the later local model, write $Z=z^3$,
$D_z=1-z^3-a^2z^6$, and
$N_g(b)=16a^4b^2-A(b-1)^4$.  The denominator-cleared value polynomial is
\[
P(b)=16D_z^2A^2b^4-\delta_\tau a^4Z^4N_g(b)^2-32A^3s^2D_z^2b^5 .
\]
% src: THEOREMS.md:729-731,975-977
Finally define the tied target value
\[
M_\tau=16-\delta_\tau c^2-16ABs^2.
\]
% src: THEOREMS.md:527-535

\textbf{PROVED: soundness.}\quad\textbf{CONDITIONAL ON SCHINZEL H:
completeness.}
Soundness, $F(z)\Rightarrow z\in\bigcup_{w\notin S}\mathfrak m_w$, is
proved.  The converse is open unconditionally; Sections~\ref{sec:record}
--\ref{sec:irreducibility} prove it assuming classical Schinzel~H.
% src: THEOREMS.md:437-439,468-475
% src: data/l22_elimination.jsonl; data/l18_schinzel_implies_h.jsonl
If that lemma holds, the union outside $S$ is $\exists_6$; adjoining the
finitely many omitted maximal ideals preserves that count, so
$\Q\setminus\Z$ is $\exists_6$, $\Z$ is $\forall_6$, and
$\operatorname{efd}_\Q(\Q\setminus\Z)\leq5$.
% src: THEOREMS.md:582-591

\subsection{Local structure and the exact global obstruction}

\textbf{PROVED (W0).}

The first criterion, W0, is uniform even in the residue field $\F_3$.
For every finite field $k$ of odd order,
\[
\sum_{a\in k}\chi(1+4a^2)=-1.
\]
% src: THEOREMS.md:477-484
Hence some $a$ makes $A$ a nonsquare, and the canonical selector is
\[
\tau_*(a)=
\begin{cases}
0,&\chi(-1)=-1,\\[2pt]
2a/A,&\chi(-1)=1.
\end{cases}
\]
% src: THEOREMS.md:484-492
It satisfies $\delta_{\tau_*}\neq0$ and
$\chi(-\delta_{\tau_*}A)=1$: the two values of
$-\delta_{\tau_*}A$ are $-A$ and $-1$.  Since $s\mapsto1+2s$ is a
bijection, tying loses no residue class and removes L5's $\F_3$ vacancy.
% src: THEOREMS.md:493-501

\textbf{PROVED (W1).}

For W1, let $w$ be odd with
$v_w(A)=v_w(\delta_\tau)=0$, $v_w(b)=1$, $v_w(z)\geq1$, and
$v_w(s)\geq0$.  The cube pullback gives either $c=0$ or
\[
v_w(c)=2v_w(a)+6v_w(z)-2\geq6v_w(z)-2\geq4.
\]
% src: THEOREMS.md:504-508
Then the tied conic
\[
-\delta_\tau A\,y^2-16B\,r^2=16-\delta_\tau c^2-16ABs^2
\]
% src: THEOREMS.md:508
is soluble over $\Q_w$ if and only if
$\chi_w(-\bar\delta_\tau\bar A)=1$.  The right side is a unit congruent
to $16$; its two subtracted terms have valuations at least $8$ and $1$.
Thus the target-place test is independent of the tied value of $s$.
% src: THEOREMS.md:509-525

\textbf{PROVED (W2).}

W2 identifies the remaining obstruction.  Put
$E_\tau=\Q(\sqrt{-\delta_\tau AB})$.  The conic is equivalent to
\[
N_{E_\tau/\mathbb Q}
\bigl(\delta_\tau A\,y+4r\sqrt{-\delta_\tau AB}\bigr)
=-\delta_\tau A\,M_\tau,
\qquad M_\tau=16-\delta_\tau c^2-16ABs^2.
\]
% src: THEOREMS.md:527-538
Its global solubility is exactly
\[
M_\tau=0\quad\text{or}\quad
-\delta_\tau A M_\tau\in N_{E_\tau/\mathbb Q}(E_\tau^\times).
\]
% src: THEOREMS.md:539-545
The zero alternative is genuine.  In the two branches the data are
\[
\begin{array}{c|c|c}
\tau& E_\tau&M_\tau\\ \hline
0&\mathbb Q(\sqrt{-AB})&16-c^2-16ABs^2\\
2a/A&\mathbb Q(\sqrt{-B})&16-c^2/A-16ABs^2.
\end{array}
\]
% src: THEOREMS.md:546-552

\subsection{Structural corollaries and the remaining lemma}

\textbf{PROVED (L7a).}
At an odd place with $v(\delta_\tau A)=v(B)=0$ and $M_\tau\neq0$, the tied
conic is soluble when $v(M_\tau)$ is even; when that valuation is odd, it
is soluble exactly when $\chi_v(-\delta_\tau AB)=1$.  If also $v(A)=0$,
the corresponding untied ternary form is universal over $\Q_v$.
% src: THEOREMS.md:602-620
Accordingly the full candidate obstruction set is
\[
\{2,\infty\}\ \cup\ \{v:\ v(M_\tau)\ \text{odd}\}\ \cup\ \{v:\ v(A)\ne0
\ \text{or}\ v(B)\ne0\}.
\]
% src: THEOREMS.md:622-645
Coefficient-bad places are genuine and cannot be replaced merely by the
support of $AB$, because valuations can cancel in that product.
% src: THEOREMS.md:622-643

\textbf{CONDITIONAL (L7b, conditional on L6 soundness).}
No fixed admissible $(s,b)$ and branch admits rational functions
$Y(z),R(z)\in\Q(z)$ satisfying
\[
-\delta_\tau A\,Y(z)^2-16B\,R(z)^2=M_\tau(z)
\]
identically.  This rules out a uniform rational section, not pointwise
coverage of all relevant rational fibres by some fixed pair.
% src: THEOREMS.md:647-671

\textbf{PROVED (L7c).}
For $M_\tau\neq0$, local solubility is equivalent to
\[
t_v:=\bigl(-16\delta_\tau AB,\ -\delta_\tau A M_\tau\bigr)_v=1,
\qquad \prod_v t_v=1.
\]
% src: THEOREMS.md:673-682
Thus a failing tied conic is obstructed at an even number of places, at
least two, and repairs must change places in pairs.
% src: THEOREMS.md:676-684

\textbf{CONDITIONAL ON SCHINZEL H / OPEN UNCONDITIONALLY (assembly).}
For every odd target $w$ and every $z$ with $v_w(z)\geq1$, one must choose
$(s,b)$ with $(1+2s,b)\in\Phi_1^{\{2\}}$,
$s\equiv(\bar a_w-1)/2\pmod{\mathfrak m_w}$, $v_w(b)=1$, and a retained
branch such that
$M_\tau=0$ or
$-\delta_\tau A M_\tau\in N_{E_\tau/\Q}(E_\tau^\times)$.
% src: THEOREMS.md:557-565
In this local formulation the free $b$-direction must steer
$E_\tau=\Q(\sqrt{-\delta_\tau AB})$ against the moving odd-valuation and
coefficient-bad places.  Later sections implement the required choice on
the fixed canonical branch: L22 supplies irreducibility, L20 supplies the
admissible class, and classical Schinzel~H supplies the member.  Thus the
weak/norm-approximation problem remains open only unconditionally, not as
an extra premise of the conditional theorem.
% src: data/l22_elimination.jsonl; data/l20_admissible.jsonl
% src: data/l18_schinzel_implies_h.jsonl
% src: THEOREMS.md:701-714

\textbf{EVIDENCE.}
Bounded witness-switching searches support this steering mechanism, but
are not a proof of assembly and are not used to upgrade the conditional
six-quantifier conclusion.
% src: THEOREMS.md:686-700
